% !TeX root = ../抽象代数笔记.tex
% 由主文档通过 \input 载入；请编译主文档。
\setcounter{chapter}{20} % 保留原书章号。
\chapter{主理想整环上的有限生成模与标准形}\label{ch:pid-modules}

\begin{description}
  \item[学习目标] 掌握有限生成模结构定理，将有限生成阿贝尔群和矩阵标准形作为同一个分解原理的应用。
  \item[先修知识] 第\ref{ch:factorization}、\ref{ch:polynomials}、\ref{ch:modules}章，以及矩阵相似与特征多项式。
\end{description}

\section{有限生成模的表示与史密斯标准形}
设 $R$ 为主理想整环，先证明有限秩自由模的子模仍为有限秩自由模，再用生成元与关系给出有限生成模的矩阵表示。通过可逆行、列变换建立史密斯标准形（Smith normal form），要求非零对角元逐项整除。说明这类矩阵等价变换不同于研究线性变换时的相似变换；在 $\mathbb Z$ 与 $F[x]$ 上给出欧几里得算法驱动的计算过程。

\section{有限生成模结构定理}
设 $M$ 为主理想整环 $R$ 上的有限生成模，证明存在 $r\geq0$ 及非零非单位 $d_1,\ldots,d_s$，满足 $d_1\mid\cdots\mid d_s$，使
\begin{equation}\label{eq:pid-module-structure}
  M\cong R^r\oplus R/(d_1)\oplus\cdots\oplus R/(d_s).
\end{equation}
允许 $s=0$；证明自由秩 $r$ 唯一，非平凡循环因子的不变因子在相伴意义下唯一。再转化为素元幂对应的初等因子分解，指出有限生成和主理想整环假设分别在哪些步骤使用。

\section{有限生成阿贝尔群的分类}
取 $R=\mathbb Z$，由式\eqref{eq:pid-module-structure}得到任意有限生成阿贝尔群为自由阿贝尔群与有限循环群的直和。回访第\ref{ch:abelian}章，将有限情形的初等证明与模论证明对照。通过关系矩阵计算商群，区分无挠部分的秩与有限挠部分的阶。

\section{有理标准形与若尔当标准形}
设 $V$ 为域 $F$ 上的有限维向量空间、$T:V\to V$ 为线性变换。证明对应的 $F[x]$-模为有限生成挠模，由不变因子得到有理标准形（rational canonical form），并读出最小多项式和特征多项式。只有当特征多项式在 $F$ 上分裂时，才在 $F$ 上进一步得到若尔当标准形（Jordan canonical form）；否则需保留有理标准形或扩张基域。

\section{例题、练习与自检}
由一个整数关系矩阵计算有限生成阿贝尔群的不变因子，再由给定最小多项式与特征多项式列出可能的矩阵标准形。训练从各次幂 $\ker(T-\lambda I)^k$ 的维数恢复若尔当块大小，其中 $\lambda$ 为特征值、$I$ 为恒等变换。综合题要求把“求商群”和“判断矩阵相似”分别翻译为模的结构问题。

